\documentclass[../main]{subfiles}

\begin{document}
% English abstract, followed by its translation (here, Spanish, as required
% by the UAM doctoral school). Acronyms follow the same rules as in the
% chapters: \acrlong for a single use, \acrfull at the first of several.
\chapter*{Abstract}
\addcontentsline{toc}{chapter}{Abstract}
This is the abstract of the thesis, organized in a few paragraphs of one sentence per source line.
It uses an acronym only once, so it spells it out with \verb|\acrlong|: \acrlong{dos}.
\lipsum[1]

\lipsum[2-3]

\begin{otherlanguage}{spanish}
\chapter*{Resumen}
\addcontentsline{toc}{chapter}{Resumen}
% In the translation, write the Spanish long form by hand and keep the
% English short form: "la densidad de estados (\acrshort{dos}, por sus
% siglas en inglés)". \acrfull/\acrlong would print the English long form.
Este es el resumen de la tesis, traducido al español, con la densidad de estados (\acrshort{dos}, por sus siglas en inglés) como ejemplo de acrónimo.
\lipsum[1]

\lipsum[2-3]
\end{otherlanguage}

\end{document}
